% ============================================================================
% Section 9.1 — What Is Probability?
% CCSS 7.SP.C.5
% ============================================================================
\section{What Is Probability?}

\begin{stepGoal}
\begin{itemize}
    \item Understand that probability is a number between $0$ and $1$
    \item Describe the likelihood of an event using its probability
\end{itemize}
\end{stepGoal}

\begin{stepVocabBox}
    \vocabItem{Probability}{A number from $0$ to $1$ that tells how likely an event is to happen.}
    \vocabItem{Event}{A specific outcome or group of outcomes you are interested in.}
    \vocabItem{Impossible}{An event that cannot happen --- probability is $0$.}
    \vocabItem{Certain}{An event that must happen --- probability is $1$.}
\end{stepVocabBox}

\begin{stepsCard}[How to Describe the Likelihood of an Event]
    \stepItem{List all \textbf{possible outcomes} of the situation.}
    \stepItem{Decide whether the event \textbf{can} happen, \textbf{must} happen, or \textbf{might} happen.}
    \stepItem{Assign a probability: $0$ (impossible), close to $0$ (unlikely), $0.5$ (equally likely), close to $1$ (likely), or $1$ (certain).}
\end{stepsCard}

% --- Worked Example 1 ---
\begin{stepExample}{A bag holds $4$ red marbles and $6$ blue marbles. Describe the probability of picking a green marble.}
    \stepShow{1} Possible outcomes: red or blue (no green marbles are in the bag).

    \stepShow{2} Picking green \textbf{cannot} happen because there are no green marbles.

    \stepShow{3} The event is impossible.
    \stepResult{$P(\text{green}) = 0$}
\end{stepExample}

% --- Worked Example 2 ---
\begin{stepExample}{A spinner has $3$ equal sections: red, blue, and yellow. What is the probability of landing on red?}
    \stepShow{1} Possible outcomes: red, blue, yellow --- that is $3$ outcomes total.

    \stepShow{2} Landing on red \textbf{might} happen --- it is one of the three equal sections.

    \stepShow{3} Red is $1$ out of $3$ equally likely outcomes: $P = \dfrac{1}{3} \approx 0.33$.
    \stepResult{$P(\text{red}) = \dfrac{1}{3}$ --- the event is possible but not likely.}
\end{stepExample}

\watchOut{Probability can never be less than $0$ or greater than $1$. If you get a number outside that range, recheck your work!}

% ============================================================================
% PRACTICE
% ============================================================================
\newpage
\begin{stepPractice}{What Is Probability? Practice}
    \resetProblems

    \stepPracticeHeader[step1Color]{List Possible Outcomes}
    \begin{multicols}{2}
    \prob A coin is flipped. List the possible outcomes. \answerBlank[2.5cm]
    \answer{Heads, Tails}
    \prob A standard die is rolled. How many possible outcomes? \answerBlank[2cm]
    \answer{$6$}
    \end{multicols}

    \stepPracticeHeader[step2Color]{Classify the Event}
    \prob What is the probability of rolling a $7$ on a standard die? \answerBlank[2cm]
    \answer{$0$ (impossible)}
    \prob What is the probability of rolling a number less than $7$ on a standard die? \answerBlank[2cm]
    \answer{$1$ (certain)}

    \stepPracticeHeader[step3Color]{Assign a Probability}
    \prob A bag has $3$ red and $7$ blue marbles. What is the probability of picking a red marble? \answerBlank[2cm]
    \answer{$\frac{3}{10} = 0.3$}

    \wordProblem{A weather app says the chance of rain tomorrow is $0.85$. Is rain likely, unlikely, or equally likely? Explain your reasoning.}{~}
    \answer{Likely --- $0.85$ is close to $1$, so rain is very likely to happen.}
\end{stepPractice}
